\documentclass{article}

\usepackage[utf8]{inputenc}
\usepackage[T1]{fontenc}
\usepackage{amsmath,amssymb,amsthm}
\usepackage{mathtools}
\usepackage{hyperref}
\usepackage{booktabs}

\makeatletter
\newcommand{\citep}[2][]{%
  \begingroup
  \def\@tmpkey{#2}%
  \ifx\@tmpkey\@empty\else
    (\ifx\@empty#1\@empty\else #1, \fi
     \@nameuse{bibkey@\@tmpkey})%
  \fi
  \endgroup
}
\expandafter\def\csname bibkey@lasser2026horizon\endcsname{Horizon~Aware}
\expandafter\def\csname bibkey@lasser2026exo\endcsname{Exogenous~Verification}
\expandafter\def\csname bibkey@lasser2026phase\endcsname{Phase~Redundancy}
\expandafter\def\csname bibkey@lasser2026paper8a\endcsname{Revealed~Sacrifice}
\expandafter\def\csname bibkey@lasser2026paper8b\endcsname{Need~Sufficiency}
\expandafter\def\csname bibkey@lasser2026micro\endcsname{Microfoundation}
\makeatother

\newtheorem{definition}{Definition}
\newtheorem{proposition}{Proposition}
\newtheorem{lemma}{Lemma}
\newtheorem{theorem}{Theorem}
\newtheorem{remark}{Remark}
\newtheorem{assumption}{Assumption}

\DeclareMathOperator{\vol}{vol}
\newcommand{\volL}{{\vol_{\mathrm{P}}}}
\newcommand{\volR}{{\vol_{\mathrm{R}}}}
\newcommand{\volRwin}{{\vol_{\mathrm{R}}^{[W]}}}
\newcommand{\volRlower}{{\vol_{\mathrm{R}}^{\mathrm{lower}}}}
\newcommand{\rext}{r_{\mathrm{ext}}}
\newcommand{\rsub}{r_{\mathrm{sub}}}
\newcommand{\rS}{r_{\mathrm{S}}}
\newcommand{\rW}{r_{\mathrm{W}}}
\newcommand{\rhomincross}{\rho_{\min}^{\mathrm{cross}}}
\newcommand{\betaL}{\beta^{\mathrm{lower}}}
\newcommand{\HHI}{\mathrm{HHI}}
\newcommand{\Pproxy}{P}
\newcommand{\Ttruth}{T}
\newcommand{\epsgap}{\varepsilon_{\mathrm{gap}}}
\newcommand{\epsnonres}{\varepsilon_{\mathrm{gap}}^{\mathrm{nonres}}}
\newcommand{\epsfloor}{\varepsilon_{\mathrm{floor}}^{\mathrm{res}}}
\newcommand{\epsnonresComposed}{\varepsilon_{\mathrm{gap},\mathrm{composed}}^{\mathrm{nonres}}}
\newcommand{\epsnonresChannel}[1]{\varepsilon_{\mathrm{gap},(#1)}^{\mathrm{nonres}}}
\newcommand{\epsnonresCoev}{\varepsilon_{\mathrm{gap},\mathrm{coev}}^{\mathrm{nonres}}}
\newcommand{\ResS}{\mathrm{ResS}}

\title{Lemma 1 (Intensive composition under co-evolution): Full Proof Draft (v2)}
\author{Paper 10 working draft for codex re-review}
\date{}

\begin{document}

\maketitle

\section{Goal and changes from v1}

Promote Lemma~1 of Paper~10 \S4.2 from inline proof outline to a
self-contained formal proof, suitable for appendix placement.

\textbf{Changes from v1 in response to codex review:}
\begin{itemize}
\item \textbf{Q1 (Definition 1).} Tightened to quantify over a
  fixed deployment class with stated operational parameters. Removed
  the over-strong "extensive" equivalence.
\item \textbf{Q2 (A1).} Added per-capability admissibility as an
  explicit hypothesis (Paper~9's Assumption~1), since the source
  paper's aggregate diagnostic ceiling does not bound the sup-norm
  without it.
\item \textbf{Q2 (A2).} Refined justification to use the direct
  ratio argument $\epsfloor = \volR(\ResS)/\volR(P) \leq 1$ rather
  than the weaker per-capability-rate justification.
\item \textbf{Q2 (A4) and main blocker.} Reformulated A4 as a
  \emph{new bounded-co-evolution assumption} on the deployment, not
  as a Paper~9 derived consequence. Paper~9's Composition
  Proposition (located in \texttt{lineage.tex}, not
  \texttt{goodhart.tex} as v1 incorrectly cited) establishes exact
  composition in the sequential regime but explicitly leaves
  co-evolution correction terms open. Paper~10's claim is
  conditional on this new assumption; honest scope.
\item \textbf{Q3.} Added explicit uniform bound $M \leq \bar{M}$
  in A4 to prevent hidden $|P|$-dependence in Step 3.
\item \textbf{Q5.} Rewrote §"Connection to the deployment claim"
  to correctly characterize Lemma~1's role: showing
  $h_{\mathrm{static}}(\theta) := \epsnonresComposed + \lambda
  \epsfloor$ is intensive, not dividing by $\mathrm{Lip}(g)$.
\end{itemize}

\section{Formal definition of intensivity}

\begin{definition}[Intensive in $|P|$ over deployment class $\mathcal{D}$]
\label{def:intensive}
Let $\mathcal{D}$ be a deployment class characterized by fixed
operational parameters (threat model, training-discipline regime,
verification infrastructure, governance protocol). A bound $B$
depending on the deployment state is \emph{intensive in $|P|$ over
$\mathcal{D}$} if there exists a constant $C \geq 0$ that may
depend on $\mathcal{D}$'s operational parameters but is independent
of $|P|$, such that $B \leq C$ for all valid deployment states in
$\mathcal{D}$.

\textbf{Note.} The constant $C$ depends on the deployment class
$\mathcal{D}$ (different threat models, training regimes, etc.\
yield different constants), but within a fixed $\mathcal{D}$, the
bound does not scale with $|P|$. This is the precise property the
deployment claim's capability-magnitude-independence requires.
\end{definition}

\section{Per-capability admissibility hypothesis}

The first source-paper hypothesis we inherit:

\begin{assumption}[Paper~9 per-capability admissibility]
\label{ass:admis}
For every capability $c \in P^{\mathrm{act}}$ in the operationally
active subspace, the per-capability proxy-truth gap $|T(c) - P(c)|$
is admissible in the sense of \citep[Microfoundation]{lasser2026micro}'s
Assumption~1 (per-capability admissibility): each capability's gap
contribution is bounded by its individual ceiling
($\epsnonresChannel{c} \leq \kappa_c$) and admissible per-channel
sub-shares.
\end{assumption}

\textbf{Note.} \citep[Microfoundation]{lasser2026micro}'s
Assumption~1 is the precondition under which the aggregate
diagnostic ceiling bounds the sup-norm $\epsnonres$. Without
per-capability admissibility, sup-norm and aggregate-ceiling
diverge, and the per-channel decomposition does not give a uniform
bound. We adopt Assumption~\ref{ass:admis} as part of Paper~10's
deployment-class specification.

\section{Source-paper intensivity assumptions}

We assume the following intensivity properties hold for deployments
in $\mathcal{D}$:

\begin{assumption}[A1: Per-channel non-residual gap intensivity, with admissibility]
\label{ass:A1}
Under Assumption~\ref{ass:admis}, for each channel $c \in \{1, 2, 3,
4\}$ in \citep[Microfoundation]{lasser2026micro}'s four-channel
decomposition, the per-channel non-residual gap contribution
$\epsnonresChannel{c}$ is intensive in $|P|$ over $\mathcal{D}$:
there exists $K_c \geq 0$ depending on $\mathcal{D}$'s operational
parameters but independent of $|P|$, such that
$\epsnonresChannel{c} \leq K_c$ for all valid deployment states in
$\mathcal{D}$.
\end{assumption}

\textbf{Justification.} Each channel's contribution is bounded by
its admissibility ceiling
(\citep[Microfoundation]{lasser2026micro}'s Assumption~1 per-channel
sub-share), which is a deployment parameter (set by verification
infrastructure thresholds). Per-channel admissibility ceilings do
not scale with $|P|$.

\begin{assumption}[A2: Residual floor intensivity (ratio form)]
\label{ass:A2}
The residual floor $\epsfloor =
\volR(\ResS)/\volR(\Pproxy)$ from
\citep[Microfoundation]{lasser2026micro}'s gap decomposition
satisfies $\epsfloor \leq 1$ under finite positive $\volR(\Pproxy)$.
In particular, $\epsfloor$ is intensive in $|P|$ over any deployment
class with $\volR(\Pproxy) > 0$, with constant $K_{\mathrm{floor}}
\leq 1$.
\end{assumption}

\textbf{Justification.} \citep[Microfoundation]{lasser2026micro}'s
$\epsfloor$ is the share of realized capability volume on the
residual class $\ResS$ (the structurally-unobservable subset).
Since $\ResS \subseteq \Pproxy$, the ratio $\volR(\ResS)/\volR(\Pproxy)
\in [0, 1]$. The residual class can grow with $|P|$, but the
\emph{normalized share} remains bounded by 1. This is the direct
ratio argument; codex's review noted this is cleaner than the
v1 per-capability-rate justification.

\begin{assumption}[A3: Bounded Lipschitz alignment property]
\label{ass:A3}
The alignment property $g$ has Lipschitz constant
$\mathrm{Lip}(g)$ that is bounded by $K_{\mathrm{Lip}} \geq 0$
independent of $|P|$ over $\mathcal{D}$. Equivalently, $g$ lies
in the class of welfare functionals whose sensitivity to single-
capability perturbations is uniformly bounded.
\end{assumption}

\textbf{Justification.} The intensivity of the deployment claim
requires a class restriction on $g$. Standard welfare functionals
(scalarizations of per-capability metrics, weighted sums with
bounded weights) satisfy A3 by construction. Pathological
functionals where $\mathrm{Lip}(g)$ grows with $|P|$ produce
extensive Goodhart bounds; A3 excludes such functionals from
$\mathcal{D}$. Paper~10 promotes A3 to explicit condition (C6) of
Theorem~\ref{thm:deployment_safety}.

\begin{assumption}[A4: Bounded co-evolution correction terms (new)]
\label{ass:A4}
We assume the deployment $\mathcal{D}$ operates in a regime where
co-evolution correction terms are bounded uniformly. Specifically:
there exists a constant $\bar{M} \geq 0$ depending on $\mathcal{D}$'s
operational parameters but \emph{independent of $|P|$}, such that
the maximum per-channel coupling magnitude satisfies
$M(\text{deployment}) \leq \bar{M}$ for all valid deployment states.

Under this bound, the co-evolution positive-part error term
satisfies:
\[
  (\epsnonresCoev)_+ \;\leq\; K_{\mathrm{coev}} \cdot \bar{M},
\]
where $K_{\mathrm{coev}} \geq 0$ is a structural constant from the
co-evolution dynamics.
\end{assumption}

\textbf{Justification (honest scope).}
\citep[Microfoundation]{lasser2026micro}'s Composition Proposition
(\texttt{lineage.tex}, in
the Composition Proposition section) establishes \emph{exact}
composition of the four channels in the \emph{sequential
intervention regime}. In the simultaneous co-evolution regime, the
proposition states that correction terms arise; the source paper
explicitly leaves the formal characterization of these correction
terms as open work.

A4 is therefore \emph{a new assumption added by Paper~10}, not a
derived consequence of the source paper. We assert that the
deployment class $\mathcal{D}$ operates in a sub-regime where
co-evolution correction is uniformly bounded by $\bar{M}$. This is
analogous to Assumption~SA1's role for the HHI surrogate: it is an
empirical/operational adequacy claim about the deployment, made
explicit so that operators see what they are conditioning on.

A4 is empirically testable in principle: deployments where $M$ is
observed to grow with $|P|$ violate A4 and fall outside the
deployment claim's scope.

\textbf{Defer for future work:} establishing a structural derivation
of A4 from finer-grained source-paper machinery, or refuting it for
some regime, would either tighten the deployment claim
(eliminating A4 as an assumption) or restrict the deployment class
$\mathcal{D}$ further.

\section{Statement of Lemma 1}

\begin{lemma}[Intensive composition under co-evolution]
\label{lem:1}
Let the composition rule $B$ be the static Goodhart bound:
\begin{equation}
\label{eq:composition_rule}
  B \;=\; \mathrm{Lip}(g) \cdot \bigl(\epsnonresComposed + \lambda \cdot \epsfloor\bigr),
\end{equation}
where $\epsnonresComposed$ is the composed non-residual gap under
co-evolution, $\epsfloor$ is the residual floor, $\lambda \in [0,1]$
is the residual weight, and $\mathrm{Lip}(g)$ is the Lipschitz
constant of the alignment property $g$.

Under per-capability admissibility (Assumption~\ref{ass:admis}) and
Assumptions~\ref{ass:A1}--\ref{ass:A4}, $B$ is intensive in $|P|$
over the deployment class $\mathcal{D}$: there exists a constant
$C^* \geq 0$ depending on $\mathcal{D}$'s operational parameters
but independent of $|P|$, such that $B \leq C^*$ for all valid
deployment states in $\mathcal{D}$.
\end{lemma}

\section{Proof}

\begin{proof}
We compose Assumptions~\ref{ass:admis}--\ref{ass:A4} into a uniform
bound on $B$.

\emph{Step 1: Decompose the composed non-residual gap.}
Under per-capability admissibility (Assumption~\ref{ass:admis}) and
the co-evolution regime, the four-channel composition produces:
\begin{equation}
\label{eq:eps_decomp}
  \epsnonresComposed \;\leq\; \sum_{c \in \{1,2,3,4\}}
  \epsnonresChannel{c} + (\epsnonresCoev)_+.
\end{equation}
The decomposition into per-channel sum plus a co-evolution
positive-part error term is the structural form
\citep[Microfoundation]{lasser2026micro}'s Composition Proposition
identifies. Under per-capability admissibility, the per-channel
contributions are well-defined sup-norm quantities.

\emph{Step 2: Apply A1 to bound the channel sum.}
By Assumption~\ref{ass:A1}, each $\epsnonresChannel{c} \leq K_c$
with $K_c$ depending on $\mathcal{D}$ but independent of $|P|$.
Therefore:
\begin{equation}
\label{eq:channel_sum_bound}
  \sum_{c \in \{1,2,3,4\}} \epsnonresChannel{c}
  \;\leq\; \sum_{c \in \{1,2,3,4\}} K_c
  \;=:\; K_{\mathrm{ch}}.
\end{equation}
$K_{\mathrm{ch}}$ is a finite, $|P|$-independent constant (over
$\mathcal{D}$).

\emph{Step 3: Apply A4 to bound the co-evolution error.}
By Assumption~\ref{ass:A4}, $M(\text{deployment}) \leq \bar{M}$
with $\bar{M}$ independent of $|P|$, and:
\begin{equation}
\label{eq:coev_bound}
  (\epsnonresCoev)_+ \;\leq\; K_{\mathrm{coev}} \cdot \bar{M}
  \;=:\; K'_{\mathrm{coev}}.
\end{equation}
$K'_{\mathrm{coev}}$ is independent of $|P|$ over $\mathcal{D}$
because both $K_{\mathrm{coev}}$ and $\bar{M}$ are.

\emph{Step 4: Combine into a bound on $\epsnonresComposed$.}
Substituting Equations~\ref{eq:channel_sum_bound}
and~\ref{eq:coev_bound} into Equation~\ref{eq:eps_decomp}:
\begin{equation}
\label{eq:eps_total_bound}
  \epsnonresComposed \;\leq\; K_{\mathrm{ch}} + K'_{\mathrm{coev}}
  \;=:\; K_{\mathrm{nonres}}.
\end{equation}
$K_{\mathrm{nonres}}$ is the sum of $|P|$-independent constants
over $\mathcal{D}$.

\emph{Step 5: Apply A2 to bound the residual floor.}
By Assumption~\ref{ass:A2}, $\epsfloor \leq K_{\mathrm{floor}} \leq
1$ over $\mathcal{D}$.

\emph{Step 6: Bound the composed slack term.}
Define the composed slack term as it appears in
Theorem~\ref{thm:deployment_safety}'s Layer~1 statement:
\begin{equation}
\label{eq:hstatic_def}
  h_{\mathrm{static}}(\theta) := \epsnonresComposed + \lambda \cdot \epsfloor.
\end{equation}
Combining Steps 4 and 5, with $\lambda \in [0,1]$:
\begin{equation}
\label{eq:hstatic_bound}
  h_{\mathrm{static}}(\theta) \;\leq\; K_{\mathrm{nonres}} +
  \lambda \cdot K_{\mathrm{floor}}
  \;\leq\; K_{\mathrm{nonres}} + K_{\mathrm{floor}}
  \;=:\; K_{\mathrm{slack}}.
\end{equation}
$K_{\mathrm{slack}}$ is independent of $|P|$ over $\mathcal{D}$.
This is the central intermediate result: the composed slack term
itself is intensive.

\emph{Step 7: Apply A3 to bound the Lipschitz constant.}
By Assumption~\ref{ass:A3}, $\mathrm{Lip}(g) \leq K_{\mathrm{Lip}}$
over $\mathcal{D}$.

\emph{Step 8: Compose into a uniform bound on $B$.}
Combining Equation~\ref{eq:composition_rule} with Steps 6 and 7:
\begin{equation}
\label{eq:final_bound}
  B \;=\; \mathrm{Lip}(g) \cdot h_{\mathrm{static}}(\theta)
  \;\leq\; K_{\mathrm{Lip}} \cdot K_{\mathrm{slack}}
  \;=:\; C^*.
\end{equation}
$C^*$ is the product of $|P|$-independent constants over
$\mathcal{D}$ and is therefore $|P|$-independent.

By Definition~\ref{def:intensive}, $B$ is intensive in $|P|$ over
$\mathcal{D}$. $\Box$
\end{proof}

\section{Discussion of constants}

The proof produces:
\[
  C^* \;=\; K_{\mathrm{Lip}} \cdot \bigl(K_{\mathrm{ch}} +
  K'_{\mathrm{coev}} + K_{\mathrm{floor}}\bigr).
\]
Each constant has an operational interpretation:
\begin{itemize}
\item $K_{\mathrm{Lip}}$: bounded sensitivity of the alignment
  property $g$ (deployment-class restriction; (C6) of Theorem~1).
  \emph{Source-derived from operational policy.}
\item $K_c$ for $c \in \{1,2,3,4\}$: per-channel admissibility
  ceilings (set by the deployment's verification infrastructure).
  \emph{Source-derived from
  \citep[Microfoundation]{lasser2026micro}'s Assumption~1.}
\item $K_{\mathrm{coev}}$: structural constant from the
  co-evolution dynamics. \emph{New constant; not derived from a
  source-paper proposition. See A4.}
\item $\bar{M}$: uniform bound on per-channel coupling magnitudes
  (deployment-class assumption A4). \emph{New constant.}
\item $K_{\mathrm{floor}} \leq 1$: residual share bound from A2.
  \emph{Source-derived from
  \citep[Microfoundation]{lasser2026micro}'s gap decomposition.}
\end{itemize}

In the spirit of codex's Q4 verdict: source-derived constants
inherit their justification from the source papers; newly assumed
constants ($K_{\mathrm{coev}}$, $\bar{M}$) are explicit assumptions
of Paper~10 and require operational verification (calibration
hooks in the deployment-tooling specification).

None of these scale with $|P|$ over the deployment class
$\mathcal{D}$.

\section{Connection to the deployment claim}

We rewrite this section per codex's Q5 to correctly characterize
Lemma~1's role.

Theorem~\ref{thm:deployment_safety}'s Layer~1 statement is:
\[
  |g(\Ttruth) - g(\Pproxy)| \;\leq\; \mathrm{Lip}(g) \cdot
  h_{\mathrm{static}}(\theta).
\]
where $h_{\mathrm{static}}(\theta) := \epsnonresComposed + \lambda
\cdot \epsfloor$ is the composed slack term defined in
Equation~\ref{eq:hstatic_def}.

Lemma~1's role is to establish that $h_{\mathrm{static}}(\theta)$ is
intensive in $|P|$ over $\mathcal{D}$:
\[
  h_{\mathrm{static}}(\theta) \;\leq\; K_{\mathrm{slack}}
  \quad \text{(Step 6 of the proof; Equation~\ref{eq:hstatic_bound})}.
\]

Since $\mathrm{Lip}(g)$ is also intensive over $\mathcal{D}$ by A3,
the product $\mathrm{Lip}(g) \cdot h_{\mathrm{static}}(\theta)$ is
intensive: bounded by $K_{\mathrm{Lip}} \cdot K_{\mathrm{slack}} =
C^*$ uniformly.

Layer~1's deployment-safety bound therefore holds uniformly across
$|P|$ over $\mathcal{D}$: this is the capability-magnitude-
independence property the deployment claim's title promises.

\textbf{Note.} Lemma~1 does \emph{not} divide the bound by
$\mathrm{Lip}(g)$ (which would fail at $\mathrm{Lip}(g) = 0$). It
shows the composed slack term itself is bounded, and the Lipschitz
multiplication is a separate composition step that preserves
intensivity by A3.

Subsequent lemmas (Lemma~2 Lyapunov-Goodhart bridge, Lemma~3 HHI
surrogate adequacy) refine the operational meaning of
$h_{\mathrm{static}}(\theta)$ without changing its intensivity.

\section{What this proof does and does not establish}

\textbf{Establishes:}
\begin{itemize}
\item A formal definition of intensivity in $|P|$ over a deployment
  class $\mathcal{D}$ (Definition~\ref{def:intensive}).
\item Per-capability admissibility as an explicit hypothesis
  (Assumption~\ref{ass:admis}), inherited from
  \citep[Microfoundation]{lasser2026micro}'s Assumption~1.
\item Four named assumptions (A1--A4) sufficient for the
  composition's intensivity, with each grounded either in
  source-paper machinery (A1, A2, A3) or as an explicit Paper~10
  assumption (A4, the bounded-co-evolution claim).
\item An explicit constant $C^*$ in terms of operational parameters,
  with no $|P|$-dependence over $\mathcal{D}$.
\item The correct connection to Theorem~1's Layer~1 statement
  (showing $h_{\mathrm{static}}(\theta)$ is intensive, not dividing
  the final bound by $\mathrm{Lip}(g)$).
\end{itemize}

\textbf{Does not establish (deferred):}
\begin{itemize}
\item Tightness of the constant $C^*$. The proof gives a uniform
  upper bound, not the tightest possible bound.
\item Existence of deployment configurations satisfying
  Assumptions~\ref{ass:admis}--\ref{ass:A4}. The assumptions are
  conditions on $\mathcal{D}$; verification is by the
  deployment-tooling specification (\S8 of Paper~10).
\item A structural derivation of A4 from source-paper machinery.
  \citep[Microfoundation]{lasser2026micro}'s Composition Proposition
  leaves co-evolution correction terms open; A4 is a new
  Paper~10 assumption, conditional on which the deployment claim
  binds.
\item Numerical values of $K_c, K_{\mathrm{coev}}, \bar{M},
  K_{\mathrm{floor}}, K_{\mathrm{Lip}}$. These are determined by
  the deployment's operational choices; the proof asserts their
  existence and $|P|$-independence over $\mathcal{D}$, but does not
  derive their numerical values.
\end{itemize}

\section{Specific verification questions for v2 codex review}

\begin{enumerate}
\item Is Definition~\ref{def:intensive} (intensive in $|P|$ over
  $\mathcal{D}$) the right reformulation? The deployment-class
  quantification was added in v2 per Q1.

\item Is Assumption~\ref{ass:admis} (per-capability admissibility)
  correctly stated as the precondition for A1? In particular, does
  \citep[Microfoundation]{lasser2026micro}'s Assumption~1 actually
  bound the sup-norm under admissibility, or does it require
  additional structural conditions?

\item Is A4's reformulation as a Paper~10 assumption (rather than a
  Paper~9 derived consequence) acceptable? Specifically: is the
  honest-scope framing — A4 is a new assumption Paper~10 conditions
  on, with empirical testability — defensible, or does it weaken
  the deployment claim more than it should?

\item Is the Layer~1 connection (\S"Connection to the deployment
  claim") now correctly stated? Specifically, is the role of
  Lemma~1 — establishing $h_{\mathrm{static}}(\theta)$ intensive,
  not dividing by $\mathrm{Lip}(g)$ — consistent with how
  Theorem~1's Layer~1 proof in
  \texttt{docs/paper10/sections/main\_theorem.tex} actually invokes
  Lemma~1?

\item Are the constant-discussion distinctions (source-derived vs.\
  newly-assumed) appropriate? Codex Q4 noted this distinction
  matters for operational deployment claims; v2 attempts to make it
  explicit.
\end{enumerate}

\end{document}
